\section{Preliminaries}
\label{sec-pre}

\begin{table}[t]
\caption{Notation.}
\label{tab-notation}
\small
\begin{tabular}{@{}p{2.15cm}p{5.6cm}@{}}
\toprule
Symbol & Meaning \\
\midrule
$(s,p,o)$, $\mathcal{T}$, $n$ & a triple; the set of triples; $n=|\mathcal{T}|$ \\
$\mathit{id}(a)$ & dictionary id of term $a$ \\
$x$, $\sigma$ & a B+ tree key, three ids in 24 bytes; a separator key \\
$e$ & an instance (an entity with a type) \\
$\pi$, $\ell(\pi)$ & a permutation (SPO, POS, OSP); the length of the prefix of $\pi$ a pattern binds \\
$B$, $F$ & capacity of a leaf (169 keys) and of an internal page (144 separators) \\
$h$ & height of a B+ tree \\
$P$, $Q$, $R$ & a page; its parent; its right sibling \\
$[a,b)$, $r$ & a half-open key range; the number of keys a scan returns \\
$H=(\mathcal{C},E)$ & the class hierarchy, a DAG of classes $\mathcal{C}$ and subclass edges $E$ \\
$\desc(c)$ & the classes at or below $c$, including $c$ \\
$\po(c)$, $\lo(c)$ & postorder number of $c$; smallest number in its subtree of the spanning forest \\
$I(c)$ & the set of intervals that covers the numbers of $\desc(c)$ \\
$E'(c)$ & non-tree edges whose tail is in $\desc(c)$ \\
$f(c)$ & fraction of terms in $\desc(c)$: $|\desc(c)|/N$ \\
$N$, $d$ & number of terms with vectors; vector dimension \\
$\mathbf{v}(t)$, $\mathrm{sim}$ & unit vector of $t$; cosine similarity \\
$q$, $k$ & the query term; the number of neighbors asked for \\
$\mathrm{inst}(c)$, $\mathrm{near}(q,c,k)$ & the instance query and the filtered similarity query \\
$\nu$, $\rho$ & IVF: number of lists; lists probed per query \\
$S$, $\tau$ & buffer-pool shards; threads \\
$m$ & product quantization: pieces per vector \\
$M$, $\ef$ & HNSW: links per node; search beam width \\
$G$, $D$ & a graph and its diameter (Lemma~\ref{lem-bw}) \\
\bottomrule
\end{tabular}
\end{table}

\stitle{Triples and patterns}.
An RDF graph is a set $\mathcal{T}$ of triples $(s,p,o)$ of subject,
predicate, and object, with $n=|\mathcal{T}|$.
A \emph{triple pattern} $(s,p,o)$ replaces some positions by variables,
written $\cdot$ here when the variable is not used elsewhere; its answers
are the triples of $\mathcal{T}$ that agree with it on every bound position.
A \emph{basic graph pattern} is a set of triple patterns joined on shared
variables.
Following common practice~\cite{weiss2008,neumann2010}, \sys replaces each
term $a$ by a dense integer $\mathit{id}(a)$ and stores triples as
three-integer keys.

\stitle{Class hierarchy}.
The \texttt{rdfs:subClassOf} triples define a directed graph
$H=(\mathcal{C},E)$ over classes, with an edge $(c,c')$ when $c$ is a
subclass of $c'$.
We write $\desc(c)$ for the classes at or below $c$, including $c$ itself, and
call a class \emph{below} $c$ when it is in $\desc(c)$.
Cycles are collapsed into one class, since the classes on a cycle are
equivalent.
The \emph{instances} of $c$ are the subjects $e$ with a triple
$(e,\texttt{rdf:type},c')$ for some $c'\in\desc(c)$.

\stitle{Similarity}.
Each term $t$ among $N$ terms has a unit vector $\mathbf{v}(t)\in\mathbb{R}^d$
produced by an embedding model, and the similarity of two terms is the cosine
$\mathrm{sim}(t,u)=\mathbf{v}(t)\cdot\mathbf{v}(u)$.

\begin{definition}[Queries]
\label{def-queries}
Given the data above, \sys answers
(i) a \emph{pattern query}, the answers of a basic graph pattern;
(ii) an \emph{instance query} $\mathrm{inst}(c)$, the instances of a class
$c$; and
(iii) a \emph{filtered similarity query} $\mathrm{near}(q,c,k)$, the $k$
terms $t\in\desc(c)$, $t\ne q$, that maximize $\mathrm{sim}(q,t)$.
\end{definition}

\noindent
Table~\ref{tab-notation} lists the notation used throughout.

